%%%%%%%%%%%%%%%%%%%%%%%%%%%%%%%%%%%%%%%%%
%% Results section
\section{Empirical Results}
\label{sec:results}

Before explicitly modelling direct and indirect effects, I first show estimates of a familiar and conventional regression estimates for between Ed PGI and demographics.
Ed PGI is correlated with years of parents' education, where those born to more educated parents have a higher Ed PGI on average.
This reflects heritability of Ed PGI, and exhibits why it is important to include years of parents' education as an explicit control in this analysis.
Similarly, for some measures of early childhood socioeconomic status, where those with lower Ed PGIs are more likely to have parents who smoked, a father who worked as a manual labourer, and have moved houses during childhood.
These demographic factors are used as controls in the mediation analysis --- referred to as $\vec X_i$ in the previous section.

\begin{figure}[h!]
    \caption{Correlation Between Ed PGI and Demographic Factors}
    \centering
    \singlespacing
    %\makebox[\textwidth][c]{
        %\includegraphics[width=\textwidth]{
            %../text/sections/figures/demographic-plot.png}}
\end{figure}

Next, in \autoref{tab:compiled-reg} I show estimates a regression of the first stage ($Z \to D$) and the reduced form ($Z \to Y$) and compare these results to estimates to a standard IV approach.
These estimates are using a binary version of the Ed PGI ($Z_i = 1$ if higher than the mean, 0 otherwise), and  binary measure for having attended higher education or not; these transformations make these estimates comparable to the mediation analysis to follow. 
The IV estimates assume that Ed PGI satisfies the exclusion restiction, a dubious assumption in this setting, and I include the IV estimates for comparison purpose.
Ed PGI strongly predicts university attendance; a high Ed PGI $Z_i = 1$ leads a 9.5\% increase in rate of having attended university.
This is after controlling for all the mentioned demographics, and mentioned other gene scores, and an F statistic of 72 shows that the Ed PGI is strong for higher education.
Similarly, Ed PGI strongly predicts annual earnings; a high score Ed PGI $Z_i = 1$ leads a 9.5\% increase in annual earnings.\footnote{
    \autoref{tab:reducedform-reg} shows the corresponding table for how reduced form estimates are moderated by the demographic and other genetic controls, which shows similar modification to the first-stage.
}
The IV estimates, using E Gene Score as an instrument for higher education, give implausibly high estimates for the returns to higher education.
This illustrates the difference between IV methods and mediation methods.

\begin{table}[h!]
    \singlespacing
    \small
    \centering
    \caption{OLS $+$ 2SLS Estimates, for Effect of Ed PGI on Education and Earnings in HRS Data.}
    %\makebox[\textwidth][c]{
\begin{tabular}{@{\extracolsep{5pt}}lccc} 
\\[-1.8ex]\hline 
\hline \\[-1.8ex] 
 & First-stage & Reduced-form & Na\"ive IV \\ 
\cline{2-4} 
\\[-1.8ex] & \multicolumn{3}{c}{Outcome:} \\ 
\\[-1.8ex] & University, $D_i$ & \multicolumn{2}{c}{$\log$ Earnings, $Y_i$} \\ 
\\[-1.8ex] & (1) & (2) & (3)\\ 
\hline \\[-1.8ex] 
EA Score, $Z_i$ & 0.095 & 0.095 &  \\ 
  & (0.014) & (0.019) &  \\ 
Attended University, $D_i$ &  &  & 1.002 \\ 
  &  &  & (0.212) \\ 
 Constant & $-$2.716 & 7.628 & 10.350 \\ 
  & (1.603) & (2.249) & (2.427) \\ 
 \hline \\[-1.8ex] 
Demographic Controls? & Yes & Yes & Yes \\ 
% Gene Score Controls? & Yes & Yes & Yes \\ 
Observations & 5,409 & 5,409 & 5,409 \\ 
R$^{2}$ & 0.157 & 0.145 & 0.044 \\ 
F Statistic (df = 14; 5394) & 71.992 & 65.382 &  \\ 
\hline 
\hline \\[-1.8ex] 
\end{tabular} 
}
    \label{tab:compiled-reg}
    \justify
    \footnotesize
    \textbf{Note}:
    This table shows OLS results of the first stage, reduced form, and 2SLS (na\"ively assuming exclusion) in the HRS analysis sample.
\end{table}

I then move to conventional mediation methods to estimate the direct and indirect channel between Ed PGI and later-life earnings, and show the estimates in \autoref{fig:mediation-plot}.
The red column shows the reduced form association between (binary) Ed PGI and log earnings; an increase from below mean to above mean Ed PGI is associated with an increase of 18\% in later-life earning, and 15\% if including demographic controls in the regression linearly or non-linearly.
Without controls, the direct genetic channel is around 9\% and the education channel 9\%.
With controls, the direct genetic channel is around 9\% and the education channel 6\%; the standard errors are large enough to rule out zero in either the direct or indirect channels, and the 95\% confidence interval of percent mediated through education is 30--50\%.

\begin{figure}[h!]
    \caption{Direct and Indirect Estimates, using OLS.}
    \centering
    \singlespacing
    %\makebox[\textwidth][c]{
        %\includegraphics[width=0.9\textwidth]{
            %../text/sections/figures/mediation-plot.png}}
    \label{fig:mediation-plot}
\end{figure}

The approach in \autoref{fig:mediation-plot} assumes sequential ignorability, although the previous section brought doubt to this assumption in observational work --- such as this setting.
As such, I perform the exact same analysis above, but instead of estimating the regressions by OLS, I use sample selection models.
The first column shows estimates of OLS with demographic controls, the second uses the same specification with a Heckman switching-selection model \citep{heckman1979sample}, and the third uses a two-stage control function estimation strategy (which is, in theory, roughly the same exercise as a sample selection model).

\begin{figure}[h!]
    \caption{Direct and Indirect Estimates, using Selection Models.}
    \centering
    \singlespacing
    %\makebox[\textwidth][c]{
        %\includegraphics[width=0.9\textwidth]{
            %../text/sections/figures/mediation-selection-plot.png}}
\end{figure}

The selection models both estimate that the direct genetic channel is indistinguishable from zero, and assign most of the association to the indirect channel (through education).
This is in contrast to the conventional mediation estimates, which estimate roughly 50-50\% between the direct and indirect channels, and show the conceptual difference in modelling unobserved confounding for education in the mediation analysis.


% \textbf{Needs a table showing how compliers differ.  See the relevant table.}
% \textbf{Result: The abadie (2003) weighting does not work for outcome $Y_i$ data, as breaking the exclusion restriction breaks independence in the proof.  But $\E[g(\vec X_i)]{D_i(1) = 1, D_i(0) = 0}$ remains identified by this approach --- just not $\E[g(Y_i)]{D_i(1) = 1, D_i(0) = 0}$}.

% \textbf{Results justifying the IV assumptions:}
% The IV assumptions (minus exclusion) are all credible in our example.
% Ed PGI significantly predicts higher education attainment (validity \ref{assumption:instrument}.\ref{assumption:validity}, see \autoref{fig:genescore-edyears});
% the Ed PGI is constructed from single--letter changes in many genes, which are plausibly random conditional on parents' education attainment (random assignment \ref{assumption:instrument}.\ref{assumption:random}, see \autoref{sec:data-mr});
% it is likely, though not independently testable, that a higher Ed PGI will not make anyone avoid higher education (monotonicity \ref{assumption:instrument}.\ref{assumption:monotonicity}).
